\chapter*{几何朗兰兹的非正式导论}
\markboth{几何朗兰兹的非正式导论}{}
\begingroup\let\thefootnote\relax\footnotetext{本章依据讲习班上的报告写成，包含三节和一个附录。}\endgroup

引入这一主题所需的全部定义，我们无法在此一一给出。读者须对代数几何、数论与表示论的基本概念有一定程度的熟悉。为弥补省略的定义，文末附有若干书目建议。

\section{函数域上 $\GL_n$ 的自守形式}

\ob{1.1}{ob:1.1}
设 $\cK$ 是正特征的整体域。也就是说，$\cK$ 是有限域 $\Fq$ 上代数曲线 $X$ 的有理函数域；这里要求 $X$ 光滑、射影、几何不可约。

记 $|X|$ 为 $X$ 的（闭）点集，它等同于 $\cK$ 的赋值。对给定点 $x$，记 $\cK_x$ 为相应的局部域，$\cO_x$ 为 $\cK_x$ 的整数环，$k_x$ 为剩余域。

回忆：$\cK$ 的阿代尔环是限制积
\[
\A:=\prod_{x\in|X|}'\cK_x .
\]

\ob{1.2}{ob:1.2}
本章只讨论代数群 $\GL_n$；相应地，$\GL_n(\A)$ 是局部紧群。粗略地说，自守形式论研究商空间 $\GL_n(\cK)\backslash\GL_n(\A)$ 上（复值）函数构成的空间。我们只考虑光滑函数，即关于某个开紧子群不变的函数，并把这空间记为
\[
\Funct(\GL_n(\cK)\backslash\GL_n(\A)).
\]
按定义，阿代尔群 $\GL_n(\A)$ 通过右平移作用在 $\Funct(\GL_n(\cK)\backslash\GL_n(\A))$ 上。

\ob{1.3}{ob:1.3}
回忆
\[
\sO:=\prod_{x\in|X|}\cO_x\subset\A
\]
称为整阿代尔环。对应的子群 $\GL_n(\sO)\subset\GL_n(\A)$ 是极大紧子群。以下固定 $\GL_n(\A)$ 上的一个 Haar 测度 $\mu$，使 $\mu(\GL_n(\sO))=1$。

按定义，$\Funct(\GL_n(\cK)\backslash\GL_n(\A))$ 中 $\GL_n(\sO)$-不变的子空间，就是双商
\[
\GL_n(\cK)\backslash\GL_n(\A)/\GL_n(\sO)
\]
上的函数空间。然而，子空间
\[
\Funct(\GL_n(\cK)\backslash\GL_n(\A)/\GL_n(\sO))\subset\Funct(\GL_n(\cK)\backslash\GL_n(\A))
\]
在 $\GL_n(\A)$ 的作用下并不稳定。但它带有 Hecke 代数 $\sH(\GL_n(\A),\GL_n(\sO))$ 的自然作用。

\ob{1.4}{ob:1.4}
先回忆局部 Hecke 代数的定义。对点 $x\in|X|$，即 $\cK$ 的一个位，Hecke 代数 $\sH(\GL_n(\cK_x),\GL_n(\cO_x))$ 是 $\GL_n(\cK_x)$ 上紧支撑的 $\GL_n(\cO_x)$-双不变函数构成的向量空间。代数结构由卷积给出：$f_1,f_2\mapsto f_1*f_2$，其中
\[
(f_1*f_2)(g)=\int_{\GL_n(\cK_x)}f_1(g\cdot g_1)\cdot f_2(g_1^{-1})\,\mu(dg_1),
\]
而 $\mu$ 是满足 $\mu(\GL_n(\cO_x))=1$ 的 Haar 测度。

由构造，$\sH(\GL_n(\cK_x),\GL_n(\cO_x))$ 是结合代数，其单位元是 $\GL_n(\cO_x)\subset\GL_n(\cK_x)$ 的特征函数。

设 $\pi$ 是群 $\GL_n(\cK_x)$（相应地 $\GL_n(\A)$）的（复）表示，$v$ 是 $\pi$ 中的向量。回忆：若 $v$ 关于某个开紧子群不变，则称 $v$ 光滑；若 $\pi$ 的所有向量都光滑，则称 $\pi$ 光滑。以下除特别声明外，$\GL_n(\cK_x)$（相应地 $\GL_n(\A)$）的表示都假定光滑。

\begin{othm}{命题}{1.5}{prop:1.5}
设 $\pi$ 是 $G(\cK_x)$ 的表示。则 $G(\cO_x)$-不动向量构成的子空间 $\pi^{G(\cO_x)}$ 自然地是 $\sH(\GL_n(\cK_x),\GL_n(\cO_x))$-模。此外，若 $\pi$ 不可约，则 $\pi^{G(\cO_x)}$ 是不可约 $\sH(\GL_n(\cK_x),\GL_n(\cO_x))$-模，或为 $0$。
\end{othm}

\begin{proof}
对 $v\in\pi$ 及 $G(\cK_x)$ 上任一局部常值紧支撑函数 $f$，可考虑作用
\[
f\cdot v=\int_{G(\cK_x)}f(g)\cdot g(v)\,\mu(g).
\]
当 $v\in\pi^{G(\cO_x)}$ 且 $f\in\sH(\GL_n(\cK_x),\GL_n(\cO_x))$ 时，$f\cdot v\in\pi^{G(\cO_x)}$，于是得到 $\sH(\GL_n(\cK_x),\GL_n(\cO_x))$ 在 $\pi^{G(\cO_x)}$ 上的作用。

现设 $\pi$ 不可约。须证每个 $v\in\pi^{G(\cO_x)}$ 都在 $\sH(\GL_n(\cK_x),\GL_n(\cO_x))$ 的作用下生成整个空间。按假设，任一 $v'\in\pi$ 都可写成 $f'\cdot v$，其中 $f'$ 是某个局部常值紧支撑函数。然而，若 $v'\in\pi^{G(\cO_x)}$，则
\[
f'\cdot v=f\cdot v,
\]
其中 $f$ 由 $f'$ 对 $G(\cO_x)$ 从左侧平均得到。
\end{proof}

\ob{1.6}{ob:1.6}
回忆：若 $\pi^{G(\cO_x)}\neq0$，则不可约 $G(\cK_x)$-表示 $\pi$ 称为球表示。容易证明，上述对应 $\pi\mapsto\pi^{G(\cO_x)}$ 给出双射：不可约球 $G(\cK_x)$-表示的同构类，与不可约 $\sH(\GL_n(\cK_x),\GL_n(\cO_x))$-模的同构类一一对应。

Gelfand 的一个基本定理与 Satake 的一个结果（对 $\GL_n$ 两者都不难）合起来，可以非常具体地描述代数 $\sH(\GL_n(\cK_x),\GL_n(\cO_x))$。首先，已知它交换。其次，它几乎就是多项式代数。确切地说，
\[
\sH(\GL_n(\cK_x),\GL_n(\cO_x))=\mathbb{C}[T_x^1,\dots,T_x^{n-1},T_x^n,(T_x^n)^{-1}],
\]
其中 $T_x^i\in\sH(\GL_n(\cK_x),\GL_n(\cO_x))$ 是 $\GL_n(\cO_x)$-双陪集的特征函数，该双陪集的对角 $n\times n$ 矩阵代表为
\[
\diag(\underbrace{\varpi_x,\dots,\varpi_x}_{i},1,\dots,1),
\]
其中第一组内 $\varpi_x$ 重复 $i$ 次，而 $\varpi_x\in\cO_x$ 是一致化子。\footnote{事实上，须把 $T_x^i$ 乘以 $(-|k_x|^{1/2})^{i(n-i)}$ 作归一化。}

于是，每个不可约 $\sH(\GL_n(\cK_x),\GL_n(\cO_x))$-模都是一维的；从而对每个球表示 $\pi$，子空间 $\pi^{G(\cO_x)}$ 也一维。

此外，不可约 $\sH(\GL_n(\cK_x),\GL_n(\cO_x))$-模的同构类，可以等同于
\[
\Spec(\mathbb{C}[T_x^1,\dots,T_x^{n-1},T_x^n,(T_x^n)^{-1}]),
\]
后者又可等同于群 $\GL_n(\mathbb{C})$ 中半单共轭类的集合。

\begin{annote}
原书脚注 1 要求把 $T_x^i$ 乘以 $(-|k_x|^{1/2})^{i(n-i)}$ 作归一化，而 \S3.2 中 Hecke 函子又带一个平移 $[i(n-i)]$。两处并非巧合：在层-函数对应下，$\ell$-进层上的 Tate 扭 $\Ql(1/2)$ 既对应权 $1/2$，又对应导出范畴中的平移 $[1]$，因此因子 $|k_x|^{i(n-i)/2}$ 恰好对应该平移。若把 \S1 的 $T_x^i$ 与 \S3.2 的 $T_x^i$ 直接对读，须记住这处归一化，否则两边的特征值相差 $q_x^{i(n-i)/2}$。\par\vspace{0.15em}{\raggedleft 【来源：原书脚注 1 与 \S3.2（A）；层-函数对应中的 Tate 扭约定（B）】\par}
\end{annote}

\ob{1.7}{ob:1.7}
现在过渡到整体 Hecke 代数，$\sH(\GL_n(\A),\GL_n(\sO))$。按定义，
\[
\sH(\GL_n(\A),\GL_n(\sO))=\bigotimes_{x\in|X|}'\sH(\GL_n(\cK_x),\GL_n(\cO_x)),
\]
其中限制张量积的意思是：考虑形如 $\bigotimes_{x\in|X|}f_x$ 的元素的有限线性组合，使得对几乎所有的 $x$，$f_x$ 是局部 Hecke 代数 $\sH(\GL_n(\cK_x),\GL_n(\cO_x))$ 中的单位元。

注意，每个不可约 $\GL_n(\A)$-表示 $\pi$ 都是限制张量积
\[
\pi\simeq\bigotimes_{x\in|X|}'\pi_x,
\]
其中几乎所有 $\pi_x$ 是球表示。同样，若 $\pi^{\GL_n(\sO)}\neq0$，则称不可约 $\GL_n(\A)$-表示 $\pi$ 为球表示；这等价于上述张量积中所有 $\pi_x$ 都是球表示。

于是，给定 $\GL_n(\A)$ 的不可约球表示 $\pi$，可以赋给它一族半单共轭类 $\{\gamma_x\}$，$x\in|X|$，$\gamma_x\in\GL_n(\mathbb{C})$。

\ob{1.8}{ob:1.8}
现在回到自守的情形，即商空间 $\GL_n(\cK)\backslash\GL_n(\A)$。回忆：上述商空间上的函数 $f$ 称为尖点的，若对所有 $g\in\GL_n(\A)$ 及所有子群 $U$ 都有
\[
\int_{U(\cK)\backslash U(\A)}f(u\cdot g)\,\mu(u)=0,
\]
其中 $U$ 取 $\GL_n$ 的标准抛物子群的幂幺根。

记尖点函数构成的子空间为 $\Funct_{\cusp}(\GL_n(\cK)\backslash\GL_n(\A))$。按定义，它关于 $\GL_n(\A)$ 的作用不变。可以证明，$\Funct_{\cusp}(\GL_n(\cK)\backslash\GL_n(\A))$ 中每个元素都自动紧支撑。

回忆：若
\[
\Hom_{\GL_n(\A)}(\pi,\Funct_{\cusp}(\GL_n(\cK)\backslash\GL_n(\A)))\neq0,
\]
则不可约表示 $\pi$ 称为尖点自守表示。

一个重要结果（多重度 $1$ 定理）说：上述 Hom 至多一维。此外，整个空间 $\Funct_{\cusp}(\GL_n(\cK)\backslash\GL_n(\A))$ 分解为不可约（从而是尖点自守）表示的直和，而不是直积分。

下文研究 $\GL_n(\A)$ 的球尖点自守表示。由上述讨论，我们的问题等价于描述空间 $\Funct_{\cusp}(\GL_n(\cK)\backslash\GL_n(\A)/\GL_n(\sO))$ 的分解——分解为算子 $T_x^i$（$i=1,\dots,n$，$x\in|X|$）的公共特征子空间之直和。我们已经知道这样的分解确实存在，且谱是单的：即每一组对应于
\[
\{\gamma_x\in\GL_n(\mathbb{C})^{ss}/\mathrm{Conj}\},\quad\forall x\in|X|
\]
的特征值，重数至多为 $1$。

另一个重要观察是：$\GL_n(\A)$ 在 $\Funct_{\cusp}(\GL_n(\cK)\backslash\GL_n(\A))$ 上的作用的调和分析，完全是代数的。这意味着，所有函数取值所在的域 $\mathbb{C}$ 的拓扑无关紧要。特别地，可以把 $\mathbb{C}$ 换成任意特征 $0$ 的域；为使特征空间分解存在，取它为代数闭域。我们的选择是 $\Ql$，其中 $(\ell,q)=1$。

\ob{1.9}{ob:1.9}
设 $\sigma$ 是 $\cK$ 的 Galois 群的 $n$ 维不可约 $\Ql$-表示。假定 $\sigma$ 处处无分歧，即它可以看作 \'etale 基本群 $\pi_1(X)\to\GL_n(\Ql)$ 的表示。按定义，这样的表示等同于 $X$ 上的 lisse $\ell$-进层（又称局部系统），记为 $E_\sigma$。

对每个点 $x\in|X|$，有确定的共轭类 $\mathrm{Fr}_x$（对应 $\pi_1(X)$ 中的 Frobenius 元）。记 $\sigma_x\in\GL_n(\Ql)^{ss}/\mathrm{Conj}$ 为 $\sigma(\mathrm{Fr}_x)$ 的半单部分。（已知：若 $\sigma$ 不可约，则族 $(\sigma_x,\,x\in|X|)$ 唯一决定 $\sigma$。）就局部系统 $E_\sigma$ 而言，$\sigma_x$ 对应 $\mathrm{Gal}(\overline{k}_x/k_x)$ 中 Frobenius 元在 $E_\sigma$ 于 $x$ 处纤维上的作用。

现在给出关键定义：设 $\sigma$ 是 $n$ 维无分歧不可约 Galois 表示，$\pi$ 是 $\GL_n(\A)$ 的尖点自守表示。若对一切 $x\in|X|$ 有
\[
\sigma_x=\gamma_x,
\]
则称 $\sigma$ 在朗兰兹的意义下对应于 $\pi$；其中 $\gamma_x\in\GL_n(\Ql)^{ss}/\mathrm{Conj}$ 是
\[
\Spec(\Ql[T_x^1,\dots,T_x^{n-1},T_x^n,(T_x^n)^{-1}])
\]
中附着于 $\pi_x$ 的相应点。

朗兰兹猜想（现已是 Lafforgue 的定理）断言：由上述对应关系决定，$\sigma$ 的集合与 $\pi$ 的集合之间存在双射。

\begin{annote}
这里把「处处无分歧」等同于「来自 $\pi_1(X)$ 的表示」，用的是曲线 $X$ 上的标准等价：$X$ 的有限 \'etale 覆盖与 $\pi_1(X)$ 的开子群一一对应，故无分歧表示就是 $\pi_1(X)$ 的表示，也就是 $X$ 上的 lisse 层；而对函数域，$\pi_1(X)$ 又可等同于 $\cK$ 的 Galois 群的最大无分歧商。之所以只取 $\sigma(\mathrm{Fr}_x)$ 的半单部分，是因为 $\mathrm{Fr}_x$ 本身只定义到共轭，且 $\sigma(\mathrm{Fr}_x)$ 的幂幺部分依赖 $x$ 的选取；只有半单部分是共轭类的不变量，才能作 Satake 参数去匹配 $\gamma_x$。\par\vspace{0.15em}{\raggedleft 【来源：Milne, \emph{\'Etale Cohomology}（A）；半单部分的动机见 \S1.6 与 \S3.4（I）】\par}
\end{annote}
